\documentclass[11pt]{article}
\usepackage[margin=1in]{geometry}
\usepackage[T1]{fontenc}
\usepackage[utf8]{inputenc}
\usepackage{amsmath,amssymb,amsthm}
\usepackage{enumitem}

\title{ICPC World Finals 2020\\A. Cardiology}
\author{}
\date{}

\begin{document}
\maketitle

\section*{Problem Summary}

Fix the pickup position $p$ of the indicated column. If the chosen card is currently in row $x$ and column
$y$, then after collecting columns and redealing rows, its new location depends only on $x$.
We must choose $p$ so that every card converges to one unique stable cell, and among all such choices
we want the stable cell closest to the center. We also need the worst-case number of iterations.

\section*{Key Observations}

\begin{itemize}[leftmargin=*]
    \item Let rows be $0$-based. If the indicated column is picked up in position $p$, then the selected card
    becomes the $((p-1)r + x)$-th card in the face-up deck, regardless of its old column.
    \item Therefore the next row is
    \[
        f(x) = \left\lfloor \frac{(p-1)r + x}{c} \right\rfloor .
    \]
    After one iteration, only the row matters.
    \item The map $f$ is monotone nondecreasing, so repeated application cannot create cycles of length
    greater than $1$; every row eventually reaches a fixed row.
    \item Write $a = (p-1)r$ and $q = c-1$. A fixed row $x$ satisfies
    \[
        x = \left\lfloor \frac{a + x}{c} \right\rfloor
        \iff qx \le a < qx + c.
    \]
    For $1 < p < c$, this gives a unique fixed row exactly when $a \bmod q \ne 0$, in which case
    \[
        x^* = \left\lfloor \frac{a}{q} \right\rfloor.
    \]
    The edge cases $p=1$ and $p=c$ always have the unique fixed rows $0$ and $r-1$.
    \item Once a card reaches row $x^*$, one more iteration sends it to the unique stable cell. So the final
    answer for the number of iterations is one plus the maximum number of row transitions needed to reach
    $x^*$.
\end{itemize}

\section*{Algorithm}

\begin{enumerate}[leftmargin=*]
    \item Iterate over every $p \in [1, c]$.
    \item Compute the unique fixed row and column:
    \begin{itemize}[leftmargin=*]
        \item if $p=1$, the stable cell is $(1,1)$;
        \item if $p=c$, the stable cell is $(r,c)$;
        \item otherwise, skip this $p$ when $((p-1)r) \bmod (c-1)=0$;
        \item for the remaining values,
        \[
            x^* = \left\lfloor \frac{(p-1)r}{c-1} \right\rfloor,\qquad
            y^* = ((p-1)r \bmod (c-1)) + 1.
        \]
    \end{itemize}
    \item Compute the distance from $(x^*+1,y^*)$ to the nearest central row and nearest central column.
    \item To find the worst-case number of iterations, simulate the row map $f$ only for the bottommost
    and topmost rows. Because $f$ is monotone, after $t$ row transitions every possible row lies inside
    the interval $[f^{(t)}(0), f^{(t)}(r-1)]$. Once both endpoints equal $x^*$, every row has reached $x^*$.
    \item The total number of card iterations is that row count plus one.
    \item Choose the candidate with minimum distance; break ties by the smallest $p$.
\end{enumerate}

\section*{Correctness Proof}

We prove that the algorithm returns the correct answer.

\paragraph{Lemma 1.}
For a fixed pickup position $p$, the next location of a card depends only on its current row.

\paragraph{Proof.}
Suppose the card is in row $x$ and in the column indicated by the spectator.
When the columns are collected, the indicated column is picked up in position $p$, so exactly $(p-1)r$
cards are placed before it. Inside that column, the card is the $(x+1)$-st card from the top.
Hence its new deck index is $(p-1)r + x$, independent of its old column. Redealing row by row makes the
next position depend only on this index, and therefore only on $x$. \qed

\paragraph{Lemma 2.}
For $1 < p < c$, there is a unique fixed row iff $((p-1)r) \bmod (c-1) \ne 0$. In that case the fixed row is
$\left\lfloor \frac{(p-1)r}{c-1} \right\rfloor$. The cases $p=1$ and $p=c$ have the unique fixed rows $0$ and
$r-1$, respectively.

\paragraph{Proof.}
Let $a=(p-1)r$ and let $x$ be a fixed row. Then
\[
    x=\left\lfloor \frac{a+x}{c}\right\rfloor
    \iff (c-1)x \le a < (c-1)x + c.
\]
For $1<p<c$, write $q=c-1$. If $a=qk+t$ with $1 \le t \le q-1$, then only $x=k$ satisfies the inequality,
so the fixed row is unique and equals $\lfloor a/q \rfloor$.
If $a=qk$ with $1 \le k \le r-1$, then both $x=k-1$ and $x=k$ satisfy the inequality, so the fixed row is
not unique.
For $p=1$, we have $a=0$, so $x=0$ is the only valid fixed row. For $p=c$, we have $a=(c-1)r$, so
$x=r-1$ is the only valid fixed row in the range $[0,r-1]$. \qed

\paragraph{Lemma 3.}
If the fixed row is unique and equal to $x^*$, then after $t$ row transitions every possible row lies in
$[f^{(t)}(0), f^{(t)}(r-1)]$.

\paragraph{Proof.}
The map $f$ is monotone nondecreasing. Therefore for every row $x \in [0,r-1]$,
\[
    f^{(t)}(0) \le f^{(t)}(x) \le f^{(t)}(r-1)
\]
for every $t \ge 0$. \qed

\paragraph{Lemma 4.}
If both endpoints $0$ and $r-1$ reach $x^*$ after $T$ row transitions, then every card reaches the stable
cell after at most $T+1$ full iterations, and some card needs exactly $T+1$ iterations.

\paragraph{Proof.}
By Lemma 3, after $T$ row transitions every row has become $x^*$. By Lemma 1, the next full iteration
sends every such card to the same stable cell determined by row $x^*$, so every card is stable after at
most $T+1$ iterations.

Conversely, because $f$ is monotone, one of the endpoint rows is the last row to reach $x^*$. Any card in
that row but not already in the stable column needs one extra iteration after its row first becomes $x^*$,
so $T+1$ iterations are sometimes necessary. \qed

\paragraph{Theorem.}
The algorithm outputs the correct answer for every valid input.

\paragraph{Proof.}
For each candidate $p$, Lemma 2 exactly characterizes whether there is a unique stable row and gives its
location. Lemma 1 gives the corresponding stable cell. Lemma 4 shows that the algorithm computes the
correct worst-case number of iterations for that $p$. The algorithm then compares all valid $p$ by the
required distance-to-center rule and breaks ties by the smallest $p$, so the reported quadruple is exactly
the one requested. \qed

\section*{Complexity Analysis}

We test all $c$ pickup positions. For each valid one we iterate the row map until the two endpoint rows
meet the fixed row; this takes $O(\log_c r + 1)$ steps. Therefore the total running time is
$O(c(\log r + 1))$, which is easily fast enough for $r,c \le 10^6$. The memory usage is $O(1)$.

\section*{Implementation Notes}

\begin{itemize}[leftmargin=*]
    \item Use 64-bit integers because $(p-1)r$ can be as large as $10^{12}$.
    \item The proof is easiest with $0$-based rows, but the output row is converted back to $1$-based.
    \item The nearest central distance is the sum of the nearest central-row distance and nearest
    central-column distance.
\end{itemize}

\end{document}
